\section{Perturbed dynamics}
\label{sec:perturbations}

\subsection{Purpose of the experiment}

Sections~\ref{sec:transfer} and~\ref{sec:proximity} live in a pure two-body
field. This part asks a specific operational question: if the manoeuvre plan is
designed in that field and flown against a more complete one \emph{without
retargeting}, how far off does it arrive? The answer sizes the mid-course
correction that a real flight plan must carry.

\subsection{Cowell formulation}

Perturbations are added directly to the acceleration in Cartesian coordinates,

\begin{equation}
\ddot{\mathbf r} = -\mu_M\frac{\mathbf r}{r^3}
                 + \mathbf a_{J_2}
                 + \mathbf a_{3B}.
\label{eq:cowell}
\end{equation}

The oblateness term is the standard second zonal harmonic \cite{vallado},

\begin{equation}
\mathbf a_{J_2} = \frac{3 J_2 \mu_M R_M^2}{2 r^4}
\begin{bmatrix}
\dfrac{x}{r}\left(5\dfrac{z^2}{r^2}-1\right)\\[8pt]
\dfrac{y}{r}\left(5\dfrac{z^2}{r^2}-1\right)\\[8pt]
\dfrac{z}{r}\left(5\dfrac{z^2}{r^2}-3\right)
\end{bmatrix}.
\label{eq:j2}
\end{equation}

The full three-component expression is retained even though the orbits are
equatorial. At $z = 0$ it collapses correctly to a purely radial correction, so
special-casing it would add a branch and remove a check.

The Earth is included as a third body in Battin's difference form \cite{battin},

\begin{equation}
\mathbf a_{3B} = \mu_E\left(\frac{\mathbf r_E - \mathbf r}{\lVert \mathbf r_E - \mathbf r\rVert^3}
                          - \frac{\mathbf r_E}{\lVert \mathbf r_E \rVert^3}\right),
\label{eq:thirdbody}
\end{equation}

with the Earth on a circular path about the Moon at the synodic rate
$n_{EM} = \sqrt{(\mu_M+\mu_E)/d_{EM}^3}$. The second term is the indirect part:
it removes the acceleration of the origin, because the frame is Moon-centred and
the Moon is itself falling towards the Earth. Dropping it is a classic error
that inflates the computed effect by orders of magnitude; retaining it is what
makes the $116\unit{m}$ result of Section~\ref{sec:perturbations:results}
meaningful.

\subsection{Acceleration budget}

Before integrating anything it is worth knowing what matters.
Table~\ref{tab:accel} evaluates the characteristic magnitudes at the chaser's
orbit radius.

\begin{table}[htbp]
\centering
\caption{Characteristic accelerations at $r = R_1 = 1837.4\unit{km}$. Oblateness
dominates the perturbation budget by more than an order of magnitude over the
Earth, which in turn dominates radiation pressure by two.}
\label{tab:accel}
\small
\begin{tabular}{lrr}
\toprule
Source & Magnitude [m/s$^2$] & Ratio to central \\
\midrule
Central lunar attraction     & \mAkep  & $1$ \\
Lunar $J_2$                  & \mAJtwo & $2.7\times10^{-4}$ \\
Earth third body             & \mATB   & $1.8\times10^{-5}$ \\
Solar radiation pressure (estimated) & $\sim 10^{-7}$ & $\sim 7\times10^{-8}$ \\
\bottomrule
\end{tabular}
\end{table}

Radiation pressure is estimated rather than integrated. For an area-to-mass
ratio of order $0.01\unit{m^2/kg}$ and a reflectivity near $1.3$, the cannonball
acceleration at $1\unit{AU}$ is around $10^{-7}\unit{m/s^2}$. It is also
intermittent: a $100\unit{km}$ lunar orbit has a period near two hours and spends
roughly $30$--$40\%$ of each revolution in shadow, so the perturbation switches
on and off rather than accumulating smoothly. Integrated over the
$\mTmissionH\unit{h}$ mission the velocity change is of order
$10^{-3}\unit{m/s}$, three orders below the correction $J_2$ already demands.
The code carries a disabled flag for it and the omission is deliberate.

\subsection{Why an equatorial orbit still feels $J_2$}

For $z \equiv 0$ equation~\eqref{eq:j2} reduces to
$\mathbf a_{J_2} = -\tfrac{3}{2}J_2\mu_M R_M^2 r^{-5}\,[x, y, 0]^\top$: a purely
radial correction. It does not tilt the plane and it does not regress the node.
What it does is stiffen the central field by a \emph{radius-dependent} amount,

\begin{equation}
\frac{\delta a_r}{a_r} = \frac{3}{2} J_2 \left(\frac{R_M}{r}\right)^2,
\label{eq:j2ratio}
\end{equation}

which is $2.73\times10^{-4}$ at the chaser radius and $2.01\times10^{-4}$ at the
target radius, so the two vehicles accumulate \emph{different} along-track phase.
Over ten and a half hours that difference is kilometres, and the miss is
therefore almost purely along-track, which is what the simulation produces.

The size of that phase drift should be taken from the secular element rates
rather than from a stiffened-central-field argument. With
$\bar n = \sqrt{\mu_M/a^3}$, $p = a(1-e^2)$ and $k = \tfrac34 \bar n J_2 (R_M/p)^2$,
the classical first-order rates are \cite{vallado}

\begin{equation}
\dot\Omega = -2k\cos i, \qquad
\dot\omega = k\,(4 - 5\sin^2 i), \qquad
\dot M = \bar n + k\sqrt{1-e^2}\,(2 - 3\sin^2 i).
\label{eq:secular}
\end{equation}

On an equatorial circular orbit the node and the argument of periapsis are
individually undefined, but their sum with the mean anomaly is not. The
observable is the mean-longitude rate

\begin{equation}
\dot\lambda = \dot\Omega + \dot\omega + \dot M
            = \bar n\left[1 + 3 J_2 \left(\frac{R_M}{a}\right)^{\!2}\right],
\label{eq:lambdadot}
\end{equation}

an excess of $\mJtwoExcess$ at the chaser radius. Comparing~\eqref{eq:lambdadot}
with the slope of $\mathrm{unwrap}\,\arctan(y/x)$ taken from the $J_2$-only
Cowell arc gives agreement to $\mJtwoRateRel$ relative, which is the check
recorded in the audit.

It is worth noting what this replaces. An earlier draft argued that a stiffer
central field shifts the mean motion by $\delta n/n = \tfrac12\,\delta\mu_{\text{eff}}/\mu$,
giving an excess of $\tfrac34 J_2 (R_M/a)^2$. That captures only the radial
stiffening and drops the periapsis and node contributions, and it is a factor of
four too small. The numerical arc agrees with~\eqref{eq:lambdadot}, not with the
hand-wave.

\subsection{Results}
\label{sec:perturbations:results}

Both vehicles are propagated over the identical Part~1 timeline with the
identical perturbed dynamics; only their initial states differ. The departure
impulse keeps its Keplerian magnitude but is applied along the \emph{perturbed}
velocity unit vector sampled at $t_{\text{wait}}$, which is what an open-loop
vehicle would actually do. No retargeting is performed.

\begin{table}[htbp]
\centering
\caption{Rendezvous miss at the nominal arrival epoch, flying the Keplerian plan
against progressively more complete dynamics. The Keplerian control case through
the identical code path closes to $\mMissTwoBody\unit{m}$, so the residuals
below are physics and not integrator noise. That control number is a
vehicle-to-vehicle separation, not the $\mKeplerErrMm\unit{mm}$ propagator
residual of \S\ref{sec:transfer:verif}; the two are named separately in
\texttt{metrics.json}.}
\label{tab:miss}
\small
\begin{tabular}{lrl}
\toprule
Dynamics & Miss & Unit \\
\midrule
Two-body control case (\texttt{miss\_LM\_MS\_twobody}) & \mMissTwoBody & m \\
Lunar $J_2$                          & \mMissJtwo     & km \\
Lunar $J_2$ + Earth third body       & \mMissJtwoTB   & km \\
\midrule
Third-body contribution alone        & \mMissTB       & m \\
Mid-course correction, retargeted (\S\ref{sec:pert:mcc}) & \mMccReal & m/s \\
\bottomrule
\end{tabular}
\end{table}

\noindent
The control case is the LM--MS separation through this same pipeline with the
perturbations switched off. It is $\mMissTwoBody\unit{m}$, not the
$\mKeplerErrMm\unit{mm}$ integrator residual of \S\ref{sec:transfer:verif};
see the note there.

Three readings follow. First, oblateness alone moves the meeting point by
$\mMissJtwo\unit{km}$ -- far beyond any docking corridor, and the reason a real
plan carries a correction. Second, the Earth contributes only $\mMissTB\unit{m}$
of that: at $100$--$400\unit{km}$ altitude the vehicle is deep enough in the
lunar well that Earth tides are a hundred times smaller than lunar oblateness.
Third, the miss is cheap to remove if it is removed early: a retargeting
impulse placed halfway through the mission costs $\mMccReal\unit{m/s}$, well
under $1\%$ of the nominal transfer budget. Section~\ref{sec:pert:mcc} derives
that number and explains why the obvious estimate is an order of magnitude too
pessimistic.

\begin{figure}[htbp]
\centering
\includegraphics[width=\textwidth]{fig11_j2_vs_kepler.png}
\caption{Keplerian design (solid) against the same plan flown with lunar $J_2$
(dashed). The tracks are visually indistinguishable for ten hours; the inset
resolves the arrival neighbourhood, where the accumulated along-track phase slip
has become a $\mMissJtwo\unit{km}$ miss.}
\label{fig:j2}
\end{figure}

\subsection{Element histories}

Classical elements are extracted along the perturbed trajectories with a
degeneracy-safe conversion: the orbits are simultaneously circular and
equatorial, so both the node vector and the eccentricity vector vanish. The
implementation falls back on the true longitude, taking the node along $+x$ when
$\lVert\mathbf n\rVert$ is below a scale-aware threshold and folding the
argument of periapsis into the argument of latitude when $e < 10^{-8}$. Without
that guard the conversion returns noise for precisely the orbits this mission
uses.

\subsection{The mid-course correction, priced properly}
\label{sec:pert:mcc}

The tempting estimate is to treat the $\mMissJtwo\unit{km}$ offset as a
relative-motion problem and cancel it with a Clohessy--Wiltshire two-impulse
over a fraction of an orbit, which costs $n_2\lVert\mathbf d\rVert = \mMccNaive\unit{m/s}$.
That number is arithmetically correct and operationally wrong. Nobody flies a
$15\unit{km}$ along-track error as a proximity manoeuvre, because it is a
\emph{phase} error and phase is bought with time, not with thrust.

A $15\unit{km}$ along-track offset at $R_2$ is an angle of
$15/2137 \approx 7\unit{mrad}$. Correcting it with a small period change over
one target revolution needs

\begin{equation}
\Delta\theta = 3\pi\,\frac{\Delta a}{a}
\;\Longrightarrow\;
\frac{\Delta a}{a} \approx 7.5\times10^{-4}
\;\Longrightarrow\;
\Delta v \approx \tfrac12 n \Delta a \sim 0.6\unit{m/s},
\label{eq:phasefix}
\end{equation}

an order of magnitude below the CW figure.

\texttt{lib/astro/midcourse\_correction.m} computes the real thing rather than
the estimate: single shooting on a correcting impulse, inside the same Cowell
model that produced the miss, targeting the \emph{perturbed} mothership at the
nominal arrival epoch. The result is $\mMccReal\unit{m/s}$ with a residual miss
of $\mMccResid\unit{m}$, which is~\eqref{eq:phasefix} to within a factor of two.

\begin{table}[htbp]
\centering
\caption{Cost of the same retarget as a function of when it is applied. The
error is accumulated during the $\mTwaitH\unit{h}$ phasing coast, so a burn
placed there has hours of lever arm; a burn halfway through the $66$-minute
transfer has $33$ minutes and is barely cheaper than the naive estimate.}
\label{tab:mcc}
\small
\begin{tabular}{lrrr}
\toprule
Impulse epoch & Lever arm [h] & Total $\dv$ [m/s] & Residual [m] \\
\midrule
Mid-mission (adopted)          & \mMccLever & \mMccReal     & \mMccResid \\
Mid-transfer                   & $0.55$     & \mMccTransfer & $< 0.1$ \\
\midrule
Naive CW cancellation          & --         & \mMccNaive    & -- \\
\bottomrule
\end{tabular}
\end{table}

The lesson is not that the CW number was miscalculated. It is that the cost of a
correction scales roughly with the inverse of the remaining lever arm, so the
question "how much does the $J_2$ miss cost?" has no answer until the epoch of
the burn is fixed. Fixing it late is what makes corrections expensive.

\subsection{Mitigation, and what was deliberately not implemented}

Three standard responses exist and are worth naming even though only the third
is quantified here.

\emph{Frozen orbits.} Choosing eccentricity and argument of periapsis so that
the secular $J_2$ rates cancel removes most of the drift for the target. Under
the real lunar field, dominated by mass concentrations rather than a clean
$J_2$, this becomes a numerical search rather than a closed form
\cite{zuber_grail}, and it is out of scope here.

\emph{Co-elliptic rendezvous.} Flying the approach on orbits of equal period
rather than equal shape makes the relative geometry insensitive to a
common-mode field error.

\emph{Mid-course correction.} The direct answer, budgeted in
\S\ref{sec:pert:mcc} at $\mMccReal\unit{m/s}$ when applied mid-mission.

\subsection{Finite-burn penalty on the departure impulse}

The impulsive assumption of Section~\ref{sec:transfer} deserves a number rather
than an apology. A finite tangential burn of duration $t_b$ sweeps an arc
$\Delta\theta = n t_b$ while thrusting, so the thrust direction rotates away
from the intended increment and the effective delta-v is the vector average,
$\sin(\Delta\theta/2)/(\Delta\theta/2)$. To first order the penalty is

\begin{equation}
\Delta v_{\text{loss}} \approx \Delta v \, \frac{(n t_b)^2}{24}.
\label{eq:gravloss}
\end{equation}

For the $\mDVone\unit{m/s}$ departure impulse, with thrust-to-weight quoted
against local lunar gravity $\mu_M/R_1^2 = \mAkep\unit{m/s^2}$:

\begin{table}[htbp]
\centering
\caption{Finite-burn penalty on the departure impulse, analytic, no integration.
Even the sluggish $T/W = 0.1$ case costs less than half a metre per second,
which is why the impulsive treatment survives.}
\label{tab:gravloss}
\small
\begin{tabular}{lrrr}
\toprule
$T/W$ (local) & Burn time [s] & Arc swept [deg] & Penalty [m/s] \\
\midrule
$0.1$ & \mGlossTb & $21.2$ & \mGlossA \\
$0.3$ & $139$     & $7.1$  & \mGlossB \\
$1.0$ & $42$      & $2.1$  & \mGlossC \\
\bottomrule
\end{tabular}
\end{table}

A note on the formula, because a dimensionally wrong version of it circulates.
The expression $\tfrac12 n^2 r t_b^2$ sometimes quoted for this estimate has
units of length, not velocity: it is the free-fall displacement over the burn.
Evaluated for the $T/W = 0.1$ case it returns $126\unit{km/s}$, which is a
useful reminder to check units before believing a formula.

\subsection{Verification}

The acceleration magnitudes of Table~\ref{tab:accel} are re-derived
independently in the audit script from the raw constants and agree to better
than $2\%$ with the closed-form expectations $\mu/R_1^2$,
$\tfrac32 J_2\mu R_M^2/R_1^4$ and $2\mu_E R_1/d_{EM}^3$. A factor of $10^3$ in
the $J_2$ row would indicate a kilometre/metre unit error; none is present.

The third-body contribution of $\mMissTB\unit{m}$ agrees with an independent
implementation to $0.2\%$, which is a stringent test of both the indirect term
in~\eqref{eq:thirdbody} and the Earth ephemeris rate. The $J_2$ miss agrees to
$4\%$; Section~\ref{sec:discussion:verif} explains why a quantity built from the
difference of two nearly equal secular slips is intrinsically that sensitive,
and why the discrepancy was documented rather than tuned away.
